\documentclass[11pt]{article}

\usepackage[margin=1in]{geometry}
\usepackage{amsmath,amssymb,amsthm,mathtools}
\usepackage{enumitem}
\usepackage{hyperref}
\usepackage[nameinlink,capitalize]{cleveref}

\hypersetup{colorlinks=true,linkcolor=blue,citecolor=blue,urlcolor=blue}

\newtheorem{theorem}{Theorem}
\newtheorem{proposition}{Proposition}
\newtheorem{lemma}{Lemma}
\newtheorem{corollary}{Corollary}
\newtheorem{definition}{Definition}
\newtheorem{remark}{Remark}
\newtheorem{warning}{Warning}

\newcommand{\F}{\mathbb F}
\newcommand{\E}{\mathbb E}
\newcommand{\TV}{\mathrm{TV}}

\title{Kummer-Fourier Smoothing and the Random-Walk Bridge}
\author{Jared Wilder}
\date{2026-05-13}

\begin{document}
\maketitle

\begin{abstract}
Paper 4 isolated a random-walk smoothing bridge (KRWS) as a named input in the conditional Erd\H{o}s--Graham \#203 route.  Papers 8--11 now make that bridge finite and exact for the Bernoulli-Kummer obstruction law actually used by the CRT sieve model.  This note proves the resulting smoothing theorem: weighted prime Hecke non-concentration gives exponential \(L^2\) and total-variation flattening of the Bernoulli obstruction distribution, with collision probability \(\ell^{-r}+O_A((\log Q)^{-A})\) in the high-\(\ell\) range once PHNC is supplied.  It also proves that the analogous deterministic prime-vector walk is not controlled by linear hyperplane non-concentration: affine hyperplane concentration gives a zero spectral gap.  Thus KRWS splits into a solved Bernoulli branch and a deterministic branch requiring affine non-concentration.  The result does not prove PHNC, Kummer-BV, BVK, or Erd\H{o}s--Graham \#203.
\end{abstract}

\section{Finite Fourier conventions}

Let
\[
G=\F_\ell^r,\qquad |G|=\ell^r,
\]
and let \(u\) denote the uniform probability measure on \(G\).  For a probability measure \(\mu\) on \(G\), define
\[
\widehat\mu(A)=\sum_{x\in G}\mu(x)e_\ell(A\cdot x),
\qquad A\in G,
\]
where \(e_\ell(t)=\exp(2\pi i t/\ell)\).  Then
\[
\|\mu-u\|_2^2
=
\ell^{-r}\sum_{A\ne0}|\widehat\mu(A)|^2,
\]
and
\[
\|\mu-u\|_{\TV}
\le
\frac12\ell^{r/2}\|\mu-u\|_2.
\]
The collision probability of \(\mu\) is
\[
\operatorname{Col}(\mu)=\sum_{x\in G}\mu(x)^2
=
\ell^{-r}+\|\mu-u\|_2^2.
\]

\section{The deterministic affine trap}

The first point is negative and structural.  Linear hyperplane non-concentration does not imply smoothing for the deterministic step measure
\[
\mu_{\mathcal Q}=\frac1{|\mathcal Q|}\sum_{q\in\mathcal Q}\delta_{x_q}.
\]

\begin{proposition}[Affine hyperplane trap]\label{P:affine-trap}
Let \(r\ge2\), let \(A_0\in\F_\ell^r\setminus\{0\}\), and let \(\mu\) be the uniform measure on the affine hyperplane
\[
H_{A_0,1}=\{x\in\F_\ell^r:A_0\cdot x=1\}.
\]
Then \(\mu\) assigns mass at most \(1/\ell\) to every linear hyperplane \(\{x:A\cdot x=0\}\), but
\[
|\widehat\mu(A_0)|=1.
\]
Consequently \(\mu^{*n}\) does not converge to \(u\) for any sequence \(n\to\infty\).
\end{proposition}

\begin{proof}
If \(A\) is a nonzero scalar multiple of \(A_0\), then \(\{A\cdot x=0\}\cap H_{A_0,1}=\varnothing\).  If \(A\) is not a scalar multiple of \(A_0\), then the two independent affine-linear equations cut \(H_{A_0,1}\) in codimension one inside \(H_{A_0,1}\), hence in relative mass \(1/\ell\).  This proves the claimed linear hyperplane non-concentration.

For \(x\in H_{A_0,1}\), the character \(e_\ell(A_0\cdot x)\) is the constant \(e_\ell(1)\).  Thus \(\widehat\mu(A_0)=e_\ell(1)\), so \(|\widehat\mu(A_0)|=1\).  Since \(\widehat{\mu^{*n}}(A_0)=\widehat\mu(A_0)^n\), the same coefficient has modulus one for every \(n\), whereas the uniform measure has zero nontrivial Fourier coefficients.
\end{proof}

\section{Deterministic smoothing under affine non-concentration}

The correct deterministic finite theorem uses affine hyperplanes, not only linear kernels.

\begin{theorem}[Affine non-concentration smoothing]\label{T:affine-smoothing}
Let \(\mu\) be a probability measure on \(G=\F_\ell^r\).  Assume that for some \(0<\delta\le1/2\),
\[
\max_{A\ne0,\ t\in\F_\ell}\mu(\{x:A\cdot x=t\})\le1-\delta.
\]
Then for every \(A\ne0\),
\[
|\widehat\mu(A)|^2
\le
1-2\delta(1-\delta)
\left(1-\cos\frac{2\pi}{\ell}\right).
\]
Consequently, for every \(n\ge1\),
\[
\|\mu^{*n}-u\|_2^2
\le
(1-\ell^{-r})
\exp\left(
-2n\delta(1-\delta)
\left(1-\cos\frac{2\pi}{\ell}\right)
\right),
\]
and
\[
\|\mu^{*n}-u\|_{\TV}
\le
\frac12(\ell^r-1)^{1/2}
\exp\left(
-n\delta(1-\delta)
\left(1-\cos\frac{2\pi}{\ell}\right)
\right).
\]
\end{theorem}

\begin{proof}
Fix \(A\ne0\), and let \(Y=A\cdot X\in\F_\ell\), where \(X\) has law \(\mu\).  Let \(\rho(t)=\mathbb P(Y=t)\).  The affine non-concentration assumption gives \(\max_t\rho(t)\le1-\delta\).  Since \(0<\delta\le1/2\),
\[
\sum_t\rho(t)^2\le(1-\delta)^2+\delta^2.
\]
Let \(Y'\) be an independent copy of \(Y\).  Then
\[
\mathbb P(Y\ne Y')
=1-\sum_t\rho(t)^2
\ge2\delta(1-\delta).
\]
Also
\[
|\widehat\mu(A)|^2
=
\E e_\ell(Y-Y').
\]
The imaginary part vanishes by symmetry, and when \(Y-Y'\ne0\) the real part is at most \(\cos(2\pi/\ell)\).  Hence
\[
|\widehat\mu(A)|^2
\le
1-\mathbb P(Y\ne Y')
\left(1-\cos\frac{2\pi}{\ell}\right),
\]
which proves the spectral gap.  Plancherel applied to \(\mu^{*n}\), together with
\[
(1-x)^n\le e^{-nx},
\]
gives the \(L^2\) bound.  The total-variation estimate follows from Cauchy--Schwarz:
\[
\|\mu^{*n}-u\|_{\TV}
\le\frac12\ell^{r/2}\|\mu^{*n}-u\|_2.
\]
\end{proof}

\begin{corollary}[Deterministic KRWS from affine PHNC]\label{C:det-krws}
For \(r=2\), if \(\ell\le(\log Q)^B\), \(\delta=\ell^{-C}\), and
\[
n\delta(1-\delta)\left(1-\cos\frac{2\pi}{\ell}\right)
\ge
A\log\log Q+\log\ell+O(1),
\]
then
\[
\operatorname{Col}(\mu^{*n})
=
\ell^{-2}+O_A((\log Q)^{-A}).
\]
Using \(1-\cos(2\pi/\ell)\ge8/\ell^2\), it is enough to take
\[
n\gg_{A,B,C}\ell^{C+2}\log\log Q.
\]
\end{corollary}

\section{The Bernoulli-Kummer obstruction walk}

The deterministic affine trap is not the obstruction law used in the CRT model.  The actual finite object is the Bernoulli-Kummer obstruction walk
\[
\nu
=
*_{q\in\mathcal Q}
\bigl((1-p_q)\delta_0+p_q\delta_{x_q}\bigr),
\]
the law of
\[
Y=\sum_{q\in\mathcal Q}B_qx_q,
\]
where the \(B_q\) are independent Bernoulli variables with parameters \(p_q\).  The independence is the CRT independence used in Paper 8.

\begin{definition}[Weighted PHNC mass]\label{D:weighted-mass}
Define
\[
W_\ell(\mathcal Q)
=
\min_{A\ne0}
\sum_{q:A\cdot x_q\ne0}p_q(1-p_q).
\]
\end{definition}

\begin{theorem}[Bernoulli-Kummer smoothing]\label{T:bernoulli-smoothing}
For the Bernoulli-Kummer obstruction law \(\nu\),
\[
\|\nu-u\|_2^2
\le
(1-\ell^{-r})
\exp\left(-\frac{16}{\ell^2}W_\ell(\mathcal Q)\right),
\]
and
\[
\|\nu-u\|_{\TV}
\le
\frac12(\ell^r-1)^{1/2}
\exp\left(-\frac{8}{\ell^2}W_\ell(\mathcal Q)\right).
\]
Equivalently,
\[
\operatorname{Col}(\nu)
=
\ell^{-r}
+O\left(
\exp\left(-\frac{16}{\ell^2}W_\ell(\mathcal Q)\right)
\right),
\]
with the explicit error factor bounded by \(1-\ell^{-r}\).
\end{theorem}

\begin{proof}
For \(A\ne0\), independence gives
\[
\widehat\nu(A)
=
\prod_{q\in\mathcal Q}
\bigl((1-p_q)+p_qe_\ell(A\cdot x_q)\bigr).
\]
Therefore
\[
|\widehat\nu(A)|^2
=
\prod_{q\in\mathcal Q}
\left[
1-2p_q(1-p_q)
\left(1-\cos\frac{2\pi(A\cdot x_q)}{\ell}\right)
\right],
\]
which is the Paper 8 squared-modulus identity after normalization by \(\varphi(P)^2\).  Using \(1-y\le e^{-y}\) and the universal gap
\[
1-\cos\frac{2\pi t}{\ell}\ge\frac8{\ell^2}
\qquad(t\ne0),
\]
we obtain
\[
|\widehat\nu(A)|^2
\le
\exp\left(
-\frac{16}{\ell^2}
\sum_{q:A\cdot x_q\ne0}p_q(1-p_q)
\right)
\le
\exp\left(-\frac{16}{\ell^2}W_\ell(\mathcal Q)\right).
\]
Plancherel over the \(\ell^r-1\) nontrivial characters gives
\[
\|\nu-u\|_2^2
=
\ell^{-r}\sum_{A\ne0}|\widehat\nu(A)|^2
\le
(1-\ell^{-r})
\exp\left(-\frac{16}{\ell^2}W_\ell(\mathcal Q)\right).
\]
The total-variation and collision estimates follow from the identities in the opening section.
\end{proof}

\begin{corollary}[PHNC high-\(\ell\) consequence]\label{C:phnc-high-ell}
Assume \(0<\alpha\le p_q\le1-\alpha\) and PHNC in the form
\[
\#\{q\in\mathcal Q:A\cdot x_q\ne0\}
\ge
\ell^{-C}|\mathcal Q|
\qquad(A\ne0).
\]
Then
\[
W_\ell(\mathcal Q)
\ge
\alpha(1-\alpha)\ell^{-C}|\mathcal Q|,
\]
and therefore
\[
\|\nu-u\|_2^2
\le
(1-\ell^{-r})
\exp\left(
-16\alpha(1-\alpha)\ell^{-C-2}|\mathcal Q|
\right).
\]
If \(\ell\le(\log Q)^B\) and \(|\mathcal Q|\asymp Q/(\ell\log Q)\), the exponent is
\[
\gg_{\alpha,B,C}
\frac{Q}{\ell^{C+3}\log Q}
\ge
\frac{Q}{(\log Q)^{B(C+3)+1}}.
\]
Thus for fixed \(A_0>0\) and fixed \(r,B,C\),
\[
\operatorname{Col}(\nu)=\ell^{-r}+O_{A_0}((\log Q)^{-A_0})
\]
for all sufficiently large \(Q\).
\end{corollary}

\section{The Bernoulli-extremal obstruction}\label{S:bernoulli-extremal}

The affine counterexample of \cref{P:affine-trap} kills deterministic finite-field smoothing.  We record here that the same failure persists under \emph{Bernoulli thinning} of the prime set, ruling out a class of extremal-only proof strategies for the Bernoulli branch.

\begin{lemma}[Bernoulli shadow]\label{L:shadow}
Let \(V=\F_\ell^r\), let \(\mathcal H\) be a finite family of affine or linear hyperplanes in \(V\), and let \(\{x_i\}\) be a finite multiset in \(V\) with independent Bernoulli weights \(\xi_i\sim\mathrm{Bernoulli}(p_i)\).  Define
\[
M=\sum_i p_i,\quad X=\sum_i\xi_i,\quad
\mu(H)=\sum_{x_i\in H}p_i,\quad X(H)=\sum_{x_i\in H}\xi_i.
\]
For \(\delta\in(0,1)\), set \(L=\log(2(|\mathcal H|+1)/\delta)\) and \(U=2\sqrt{ML}+2L\).  Then with probability \(\ge 1-\delta\), simultaneously for every \(H\in\mathcal H\),
\[
|X-M|\le U,\qquad |X(H)-\mu(H)|\le U.
\]
If \(\mu(H_0)\ge(\theta+\eta)M\) for some \(H_0\) and \(U\le\eta M/6\), then \(X(H_0)/X\ge\theta+\eta/2\).
\end{lemma}

\begin{proof}
Bernstein's inequality applied to \(X(H)-\mu(H)=\sum_i\mathbf 1[x_i\in H](\xi_i-p_i)\), then a union bound over \(\mathcal H\cup\{V\}\).
\end{proof}

\begin{corollary}[Bernoulli does not repair hyperplane support]\label{C:bernoulli-no-repair}
If every \(x_i\) with \(p_i>0\) lies in a common hyperplane \(H_0\), then \(X(H_0)=X\) identically.  If \(M\to\infty\), then \(\mathbb P(X=0)\le\exp(-M)\to 0\), so the Bernoulli sample is, with probability \(1-o(1)\), nonempty and \(100\%\) hyperplane-concentrated.
\end{corollary}

We exhibit an explicit arithmetic family where this happens.

\begin{lemma}[Hyperplane prime family]\label{L:hyperplane-family}
Fix an odd prime \(\ell\) and \(A=(A_1,A_2)\in\F_\ell^2\setminus\{0\}\).  Set
\[
\alpha_A=2^{A_1}3^{A_2},\qquad K_{\ell,\star}=\mathbb Q(\zeta_\ell,2^{1/\ell},3^{1/\ell}).
\]
By rank-two Kummer independence (companion Paper~6), \(\operatorname{Gal}(K_{\ell,\star}/\mathbb Q(\zeta_\ell))\simeq\F_\ell^2\).  The prime set
\[
\mathcal P_{\ell,A}=\{q\nmid 6\ell:q\equiv 1\bmod\ell,\ x_\ell(q)\in H_A\},
\qquad H_A=\{v\in\F_\ell^2:A\cdot v=0\},
\]
has natural density \(\delta_{\ell,A}=1/(\ell(\ell-1))\) by qualitative Chebotarev for \(K_{\ell,\star}/\mathbb Q\).
\end{lemma}

\begin{proof}
The hyperplane atom criterion (Paper~6) gives \(x_\ell(q)\in H_A\Leftrightarrow\alpha_A\in(\F_q^*)^\ell\).  The absolute Galois group of \(K_{\ell,\star}/\mathbb Q\) has order \(\ell^2(\ell-1)\); the conditions \(q\equiv 1\bmod\ell\) (factor \(1/(\ell-1)\)) and \(x_\ell(q)\in H_A\) (factor \(|H_A|/|V|=1/\ell\)) compose to density \(1/(\ell(\ell-1))\).
\end{proof}

\begin{theorem}[Intrinsic Bernoulli-Kummer hyperplane obstruction]\label{T:bernoulli-obstruction}
Fix an odd prime \(\ell\) and nonzero \(A\in\F_\ell^2\).  For \(X\ge 2\) let
\[
\mathcal Q_X=\mathcal P_{\ell,A}\cap[1,X],\qquad P_X=\prod_{q\in\mathcal Q_X}q,
\]
and choose \(m\) uniformly in \((\mathbb Z/P_X\mathbb Z)^*\).  Define
\[
W_X=\sum_{q\in\mathcal Q_X}B_q(m),\qquad W_X(H_A)=\sum_{\substack{q\in\mathcal Q_X\\ x_\ell(q)\in H_A}}B_q(m).
\]
Then \(W_X(H_A)=W_X\) identically, and \(\mathbb P_m(W_X>0)\to 1\) as \(X\to\infty\).
\end{theorem}

\begin{proof}
By \cref{L:hyperplane-family}, every \(q\in\mathcal Q_X\) satisfies \(x_\ell(q)\in H_A\), so \(W_X(H_A)=W_X\) identically.  By Mertens applied to a density-\(1/(\ell(\ell-1))\) prime set,
\(M_X=\mathbb E_m W_X\ge\sum_{q\in\mathcal Q_X}1/(q-1)=(\ell(\ell-1))^{-1}\log\log X+O(1)\to\infty\),
hence \(\mathbb P(W_X=0)\le\exp(-M_X)\to 0\).
\end{proof}

\begin{corollary}[Bernoulli-extremal bypass kill]\label{C:bernoulli-kill}
No proof of Bernoulli-weighted PHNC at level \(\ell\) using only extremal arguments on the Kummer log vectors \(\{x_\ell(q)\}\subset\F_\ell^2\) can avoid producing the deterministic \(p_q\)-weighted Kummer equidistribution estimate
\[
\max_{A\ne 0}\Biggl|\sum_{\substack{q\le X\\ q\equiv 1\bmod\ell}}p_q\Bigl(\mathbf 1_{A\cdot x_\ell(q)=0}-\tfrac1\ell\Bigr)\Biggr|=o\!\!\left(\sum_{\substack{q\le X\\ q\equiv 1\bmod\ell}}p_q\right).
\]
\end{corollary}

\begin{proof}
Suppose an extremal-only proof exists.  Specialise to the prime set \(\mathcal Q_X=\mathcal P_{\ell,A_0}\cap[1,X]\) for some chosen \(A_0\).  By \cref{T:bernoulli-obstruction}, the Bernoulli sample is hyperplane-concentrated on \(H_{A_0}\) with probability \(1-o(1)\).  Bernoulli-weighted PHNC restricted to \(\mathcal Q_X\) would force this concentration to be incompatible with prime distribution, which is impossible unless the proof controls the deterministic hyperplane discrepancy on \(\mathcal Q_X\) — the displayed quantity.
\end{proof}

The takeaway is methodological: the Bernoulli branch of KRWS (\cref{T:bernoulli-smoothing}) is the correct finite reduction, but the deterministic weighted Kummer equidistribution input it requires cannot be replaced by purely extremal \(\F_\ell^2\) combinatorics, even after Bernoulli thinning.

\section{KRWS status}

Paper 4's KRWS input should be read in two branches.

\begin{enumerate}[leftmargin=2em]
\item \textbf{Bernoulli branch.}  For the obstruction vector \(V(m)=\sum_qB_q(m)x_q\), KRWS is the finite theorem \cref{T:bernoulli-smoothing}, conditional only on the same weighted PHNC input used in Paper 11.  No affine hypothesis is needed, because every local step is anchored by the \(0\)-increment option.
\item \textbf{Deterministic branch.}  For the raw prime-vector step measure \(\mu_{\mathcal Q}=|\mathcal Q|^{-1}\sum_q\delta_{x_q}\), linear hyperplane non-concentration is insufficient by \cref{P:affine-trap}.  The correct finite replacement is affine non-concentration, which gives \cref{T:affine-smoothing}.
\end{enumerate}

This closes the finite \(L^2\) collision-smoothing part of the Bernoulli-Kummer route under PHNC.  It does not close the analytic sieve-transfer part.  KRC, RTHLS, Kummer-BV, and BVK still encode the passage from finite obstruction-law smoothing to estimates for the von Mangoldt-weighted sequence \(m2^k3^\ell+1\).

\section{Dependency status}

\begin{itemize}[leftmargin=2em]
\item The affine counterexample is unconditional and finite.
\item The deterministic smoothing theorem is unconditional once affine non-concentration is assumed.
\item The Bernoulli-Kummer smoothing theorem is unconditional once the weighted PHNC mass \(W_\ell(\mathcal Q)\) is supplied.
\item The PHNC high-\(\ell\) corollary is conditional on PHNC and a density floor for the local Bernoulli parameters.
\item No unconditional proof of PHNC, Kummer-BV, BVK, or Erd\H{o}s--Graham \#203 is asserted.
\end{itemize}

\end{document}
